% concatenated from 0830.tex
\documentclass[11pt]{amsart}

\usepackage{enumerate}
\usepackage{amsmath,amssymb}
\usepackage{comment}
\usepackage{bm}
\usepackage{xcolor}
\newcommand{\oi}[1]{{\color{teal}#1}}
%\usepackage{combelow}

\newtheorem{theorem}{Theorem}
\newtheorem{lemma}[theorem]{Lemma}
\newtheorem{cor}[theorem]{Corollary}
\newtheorem{proposition}[theorem]{Proposition}
\newtheorem{problem}[theorem]{Problem}
\newtheorem{definition}[theorem]{Definition}


\theoremstyle{remark}
\newtheorem{remark}[theorem]{Remark}

%%%%%%%%%%    author's macro     %%%%%%%%%
\newcommand{\B}{\mathcal{B}}
\newcommand{\BH}{\mathcal{B(H)}}
\newcommand{\bC}{\mathbb{C}}
\newcommand{\bN}{\mathbb{N}}

\newcommand{\bR}{\mathbb{R}}
\newcommand{\FH}{\mathcal{F(H)}}
\newcommand{\cH}{{\mathcal{H}}}
\newcommand{\PH}{\mathcal{P}_1(\cH)}
\newcommand{\SH}{\mathcal{S(H)}}
\newcommand{\ADj}{\mathcal{A}_{D_j}}
\newcommand{\AC}{\mathcal{A}_C}
\newcommand{\AD}{\mathcal{A}_D}

\newcommand{\SC}{\mathcal{S}_C}
\newcommand{\SD}{\mathcal{S}_D}
\newcommand{\SE}{\mathcal{S}_E}
\newcommand{\Conj}{\operatorname{Conj}(\cH)}
\newcommand{\SCY}{\mathcal{S}_{C,1}}
\newcommand{\SCO}{\mathcal{S}_{C,0}}
\newcommand{\SCq}{\mathcal{S}_{C,q}}
\newcommand{\SSC}{\mathcal{O}_C}
\newcommand{\OC}{\mathcal{O}_C}
\newcommand{\cA}{\mathcal{A}}
\newcommand{\cJ}{\mathcal{J}}
\newcommand{\cL}{\mathcal{L}}
\newcommand{\cI}{\mathcal{I}}
\newcommand{\cBH}{\mathscr{B}_1(\cH)}
\newcommand{\dBH}{\mathscr{B}_0(\cH)}
\newcommand{\MC}{\mathcal{M}_n(\bC)}

\newcommand{\eps}{\varepsilon}
\newcommand{\E}{\mathscr{V}}
\newcommand{\U}{\mathcal{Z}}
\newcommand{\F}{\mathscr{F}}

\newcommand{\ind}{\textup{ind}}
\newcommand{\cF}{\mathcal{F}}
\newcommand{\cS}{\mathcal{S}}
\newcommand{\la}{\langle}

\newcommand{\M}{{\mathcal{M}}}
\newcommand{\onb}{\textsc{onb}}
\newcommand{\ra}{\rangle}

\newcommand{\Span}{\operatorname{span}}


\newcommand{\ran}{\textup{ran}}

\newcommand{\rank}{\textup{rank}}
\newcommand{\tr}{\textup{tr}}

\begin{document}
\title[]{Linear maps preserving $C$-symmetry operators}

\author[M.Anzai]{Mayu Anzai}
\address{Graduate School of Science and Technology, 
Niigata University, Niigata 950-2181, Japan.}
\email{f25a046a@mail.cc.niigata-u.ac.jp}


\author[S.Oi]{Shiho Oi}
\address{Department of Mathematics, Faculty of Science, 
Niigata University, Niigata 950-2181, Japan.}
\email{shiho-oi@math.sc.niigata-u.ac.jp}



\thanks{The second author was supported in part by JSPS KAKENHI Grant Number JP24K06754.
}

%\thanks{}
\subjclass[2020]{Primary  47B02, 47B49.} 

\keywords{$C$-symmetry operators, Diagonal operators, Linear preserver}

\date{}

\begin{abstract}
 Let $\cH$ be a separable complex Hilbert space.  In this paper, we study a  continuous bijective linear map $T$ on the algebra of all bounded linear operators on $\cH$.  We characterize the map $T$ that maps the set of all $C$-symmetric operators  onto the set of all $\psi(C)$-symmetric operators, where $\psi$ is a commutativity-preserving bijection on the set of all conjugations. Furthermore, we show that this condition is equivalent to the existence of a bijection $\varphi$ on the set of all orthonormal bases of $\cH$ such that $T$ maps the set of all $\{e_n\}$-diagonal operators  onto  the set of all $\varphi(\{e_n\})$-diagonal operators.
\end{abstract}

\maketitle

\section{Introduction}
Throughout this paper, we assume that $\cH$ is a complex separable Hilbert space and $\BH$ is the algebra of all bounded linear operators on $\cH$. We denote the identity operator on $\cH$ by $I$. We first require a few preliminary definitions.
\begin{definition}
A map $C:\cH \rightarrow\cH$ is called a conjugation if
\begin{enumerate}
  \item[(i)] $C$ is antilinear, i.e. $C(\alpha x+y)=\overline{\alpha} C x+ Cy$ for $x,y\in\cH$ and $\alpha\in\bC$,
  \item[(ii)] $C$ is invertible with $C^{-1}=C$, and
  \item[(iii)] $\la Cx, Cy\ra=\la y,x\ra$ for all $x, y\in\cH$.
\end{enumerate}
We denote the set of all conjugations on $\cH$ by $\Conj$. Given an orthonormal basis $\{e_n\}$, a conjugation $C$ satisfying $Ce_n=e_n$ for all $n$ is called  a conjugation associated with $\{e_n\}$.  
\end{definition}


\begin{definition}
An operator $A \in\BH$ is said to be complex symmetric if
$CAC=A^*$ for some conjugation $C$ on $\cH$; in this case, $A$ is called  $C$-symmetric. For each conjugation $C$ on $\cH$, let $\SC$ denote the set of  $C$-symmetric operators on $\cH$. 
\end{definition}

%Then $\SC$ is a (weak-operator) closed subspace of $\BH$. 
Recall that the set of all complex symmetric operators is given by $\bigcup_{C \in \Conj} \SC$. For any conjugation $C \in \Conj$ on $\cH$,  $\SC$ is a complex linear subspace of $\BH$ that is closed under the involution $*$. It is well known that $A$ is a complex symmetric operator if and only if there is an orthonormal basis of $\cH$ with respect to which $A$ has a complex symmetric matrix representation. Garcia and Putinar initiated the study of complex symmetric operators. They showed that the class of complex symmetric operators contains many important classes of operators (see \cite{GP}). %, including normal operators, Hankel operators, compressed Toeplitz operators, and some Volterra operators.

Linear preserver problems have been extensively studied over the past several decades on matrix algebras $\MC$ or $\BH$. 
In particular, linear maps preserving normal operators were characterized in \cite{Brs}. 
Subsequently, linear maps preserving diagonalizable matrices were characterized in \cite{OS}. 
Recently, Amara, Oudghiri, and Souilah studied maps on $\BH$ that preserve the class of $C$-symmetric operators for every conjugation $C$, or equivalently, the class of $\{e_n\}$-diagonal operators for every orthonormal basis $\{e_n\}$ in \cite{AOS}. Recall that, given an orthonormal basis $\{e_n\}$ of $\cH$, an operator $D \in \BH$ is called $\{e_n\}$-diagonal if it has a diagonal matrix representation with respect to the basis $\{e_n\}$. We simply call $D$ is diagonal if it is $\{e_n\}$-diagonal for some orthonormal basis $\{e_n\}_{n}$ of $\cH$. We denote the set of all orthonormal bases of $\cH$ by $\mathcal{E}(\cH)$, and the set of all $\{e_n\}$-diagonal operators by $\mathcal{D}_{\{e_n\}}$ for $\{e_n\} \in \mathcal{E}(\cH)$. Here, we regard an orthonormal basis as an ordered sequence. Thus, $\{e_1, e_2, \cdots \}$ and $\{e_2, e_1, \cdots \}$ are regarded as distinct elements of $\mathcal{E}(\cH)$. 
In \cite{AOS}, they proved the following theorem. 
\begin{theorem}{\cite[Theorem 1]{AOS}}\label{AOST}
Let $\cH$ be a complex separable Hilbert space with $\dim \cH \ge 3$ and $T: \BH \to \BH$ be an additive map. Then the following are equivalent.
\begin{itemize}
\item[(1)] $T(\SC) \subset \SC, \ \text{for every} \ C \in \Conj$,
\item[(2)] $T(\mathcal{D}_{\{e_n\}}) \subset \mathcal{D}_{\{e_n\}},  \ \text{for every} \ \{e_n\} \in \mathcal{E}(\cH)$,
\item[(3)] there exist $\alpha, \beta \in \mathbb{C}$ and an additive functional $f$ on $\BH$ such that 
\[
T(A)=\alpha A+\beta A^{*}+f(A)I, \text{for all} \ A \in \BH.
\]
\end{itemize}
\end{theorem}

In this paper, we investigate what happens when the condition 
$T(\SC) \subset \SC$ for every $C \in \Conj$
is replaced by the more general condition that  
\[
T(\SC) \subset \mathcal{S}_{\psi(C)},
\]
where $\psi:\Conj \to \Conj$ is a bijection. To state our main result, we introduce the following definitions.

\begin{definition} 
We define the Pauli matrices by
\[
\sigma_1=
\begin{pmatrix}
0&1\\
1&0
\end{pmatrix},
\qquad
\sigma_2=
\begin{pmatrix}
0&-i\\
i&0
\end{pmatrix},
\qquad
\sigma_3=
\begin{pmatrix}
1&0\\
0&-1
\end{pmatrix}. 
\]
Then every matrix $A\in M_2(\mathbb{C})$ can be written uniquely in the form
\[
A=\frac{1}{2}(aI+x_1\sigma_1+x_2\sigma_2+x_3\sigma_3), \quad a, x_1, x_2, x_3 \in \mathbb{C}. 
\]
We shall write
\[
A=\frac{1}{2}(aI+x\cdot \sigma), 
\]
where  $x=(x_1, x_2, x_3) \in \mathbb{C}^3$ and $\sigma=(\sigma_1, \sigma_2, \sigma_3)$. 
\end{definition}
It is well-known that $\frac{1}{2}(aI+x\cdot \sigma)$ is a rank one projection if and only if $a=1$ and $x=(x_1, x_2, x_3) \in \mathbb{R}^3$ with $x_1^2+x_2^2+x_3^2=1$. 

We now present the main results of this paper.

\begin{theorem}\label{main0}
Let $\cH$ be a complex Hilbert space with $\dim \cH=2$.  Let $T: M_2(\mathbb{C}) \to M_2(\mathbb{C})$ be a bijective linear map. Then the following are equivalent.
\begin{itemize}
\item[(1)] there is a bijection $\psi: \Conj \to \Conj$ which satisfies  
 \[
T(\SC) = {\mathcal{S}_{\psi(C)}},
\]
\item[(2)] there is a bijection $\varphi: \mathcal{E}(\cH) \to \mathcal{E}(\cH)$,  which satisfies  
\[
T(\mathcal{D}_{\{e_n\}}) = \mathcal{D}_{\varphi(\{e_n\})},
\] 
\item[(3)] there exists $c, \lambda \in \mathbb{C} \setminus \{0\}$, a linear functional $\ell$ on $\mathbb{C}^3$, and $B \in GL(3,\mathbb{R})$ such that 

\[
T(aI+x\cdot \sigma)=c aI+\ell(x)I+\lambda Bx
\] 
for  $a \in \mathbb{C}$ and $x \in \mathbb{C}^3$.
\end{itemize}

\end{theorem}

\begin{definition}
A bijection $\psi: \Conj \to \Conj$ is said to preserve commutativity if 
\[
 C_1C_2=C_2C_1 \Longleftrightarrow \psi(C_1)\psi(C_2)=\psi(C_2)\psi(C_1)
 \]
for all $C_1, C_2 \in \Conj$.

\end{definition}

%When $\cH \ge 3$, we have the following.
\begin{theorem}\label{main}
Let $\cH$ be a complex separable Hilbert space with $\dim \cH \ge 3$.  Let $T: \BH \to \BH$ be a continuous bijective linear map. Then the following are equivalent.
\begin{itemize}
\item[(1)] there is a bijection $\psi: \Conj \to \Conj$ preserving commutativity, which satisfies  
 \[
T(\SC) = {\mathcal{S}_{\psi(C)}},
\]
\item[(2)] there is a bijection $\varphi: \mathcal{E}(\cH) \to \mathcal{E}(\cH)$,  which satisfies  
\[
T(\mathcal{D}_{\{e_n\}}) = \mathcal{D}_{\varphi(\{e_n\})},
\] 
\item[(3)] there exists $c \in \mathbb{C} \setminus \{0\}$, a bounded linear functional $f$ on $\BH$, and a unitary or anti-unitary $U \in \BH$ such that 
\[
T(A)=cUAU^{*}+f(A)I, \quad A \in \BH.
\] 
\end{itemize}
\end{theorem}
%The implications $(3)\Rightarrow(1)$ and  $(3)\Rightarrow(2)$ are obvious.
%The case $\dim \cH=1$ is trivial. Hence, we may assume that $\dim \cH \ge 2$. %In this paper, we show that $(1)\Rightarrow(2)$,  $(2)\Rightarrow(3)$ and $(3)\Rightarrow(1)$.  


For convenience, we introduce the following notation, which will be used throughout the paper. By $\sigma(A)$,  we denote the spectrum of $A$. %We denote the subspace spanned by some elements of $V$ by $\operatorname{span}V$, where $V$ is a set of operators in $\BH$. 
We denote $\mathbb T=\{\lambda\in\mathbb C \mid |\lambda|=1\}$. Set $N=\dim\cH$, and define $\mathcal{I}_N=\{1,\ldots,N\}$ if $N<\infty$,  and $\mathcal{I}_N=\mathbb{N}$ if $N=\infty$.


\section{Preliminaries}
In this section, we first present some preliminary results that hold for $\dim \cH \ge 2$.
\begin{lemma}{\cite[Lemma 1]{GP}}\label{GP1}
If $C$ is a conjugation on $\cH$, then there exists an orthonormal basis $\{e_n\}$ such that $Ce_n=e_n$ for all $n$.
\end{lemma}


\begin{lemma}\label{I}
Let $A \in \BH_{sa}$. There is $c \in \mathbb{R}$ such that $A=cI$ if and only if $A \in \SC$ for all conjugations $C$. 
\end{lemma}

\begin{proof}
We consider two cases: (i) $\dim \cH=2$ and (ii) $\dim \cH \ge 3$.\\
%Let be a conjugation $C$. By Lemma \ref{GP1}, there is an orthonormal basis $\{e_n\}$ such that $Ce_n=e_n$ for all $n$. There is a unitary such that $Ua_n=e_n$ for all $n$. We have $UJ_0U^{*}e_n=e_n$, which implies that $UJ_0U^{*}=C$. By assumption we have $CAC=A$, thus $U^{*}AU=U^*UJ_0U^{*}AUJ_0U^{*}U=J_0U^{*}AUJ_0$. 
(i) Suppose $\dim \cH=2$. Assume that $A \in \SC$ for all conjugations $C$. Let $U$ be a unitary matrix such that 
\[
U^{*}AU=\begin{pmatrix}
\lambda_1 & 0  \\
0 & \lambda_2 
\end{pmatrix},
\]
 where $\sigma(A)=\{ \lambda_1, \lambda_2\}$. Let $\{(1,0) ,(0,1)\}$ be the standard basis of $\mathbb{C}^2$. 
A conjugation $J$ on $\mathbb{C}^2$ is defined by
\[
J(z_1, z_2)
=
(\overline{z_2}, \overline{z_1})
\]
for every $(z_1,z_2)\in\mathbb{C}^2$, where $(z_1,z_2)$ denote the coordinate representation with respect to the standard basis $\{(1,0) ,(0,1)\}$.  Then we get a conjugation $UJU^{*}$.  Since $A=UJU^{*}AUJU^{*}$, we get $U^{*}AU=JU^{*}AUJ$.  This implies 
\[
\begin{pmatrix}
\lambda_1 & 0 \\
0 & \lambda_2 
\end{pmatrix}=\begin{pmatrix}
\lambda_2 & 0 \\
0 & \lambda_1 
\end{pmatrix}.
\]
Thus we get $\lambda_1=\lambda_2$, which implies $A=cI$ for some $c \in \mathbb{R}$. \\

(ii) Suppose $\dim \cH \ge 3$. We assume that  $A \in \SC$ for all conjugations $C$.  Let $x \in \cH$ with $x \neq 0$. Then there are $\alpha \in \mathbb{C}$ and $y \in (\operatorname{span} \{x\} )^{\perp} $ such that
\[
Ax= \alpha x + y.
\]
Suppose that $y \neq 0$. As $\dim \cH \ge 3$, there is $z \in (\operatorname{span} \{x, y\} )^{\perp} $ with $\|z\|=1$. 
Let $C$ be a conjugation such that $Cx=x$ and $C(\frac{y}{\|y\|})=z$. We have 
\[
CAx=ACx=Ax.
\]
Thus 
\[
\overline{\alpha} x + \|y\|z= C(Ax)=\alpha x + y, 
\]
which is a contradiction. Thus $y=0$. Since $A \in \BH_{sa}$, we obtain $\alpha \in \mathbb{R}$. 
We define $\alpha(x) \in \mathbb{R}$ for any $x \in \cH \setminus \{0\}$  by 
\[
Ax= \alpha(x)x.
\]
Let $x, y \in \cH \setminus \{0\}$. If $\dim \operatorname{span} \{x,y\}=1$, then there is $\lambda \in \mathbb{C}$ such that $x=\lambda y$. We get 
\[
\alpha(x)x=A(x)=A(\lambda y)=\lambda A(y)=\lambda \alpha(y) y=\alpha(y) x.
\]
As $x \neq 0$, we obtain $\alpha(x)=\alpha(y)$. If $\dim \operatorname{span} \{x,y\}=2$, there is a non-zero vector $z \in \operatorname{span} \{x,y\} ^{\perp} \setminus \{0\}$ as $\dim \cH \ge 3$. Then there are orthogonal bases of $\cH$, $\{x,z, \cdots\}$ and $\{y, z, \cdots\}$. We have
\[
\alpha(x+z)(x+z)=A(x+z)=A(x)+A(z)=\alpha(x)x+\alpha(z)z.
\] 
It yields that 
\[
\alpha(x)=\alpha(x+z)=\alpha(z).
\]
By a similar argument for $\{y, z, \cdots\}$, we have
\[
\alpha(y)=\alpha(y+z)=\alpha(z).
\]
Thus we get $\alpha(x)=\alpha(z)=\alpha(y)$. As a conclusion, we have $\alpha(x)=\alpha(y)$ for any $x, y \in \cH \setminus \{0\} $. Thus $\alpha=\alpha(x)$ for any $x \in \cH \setminus \{0\}$.  Thus $A=\alpha I$. The converse follows immediately.
\end{proof}



\begin{lemma}\label{dim1}
Let $\Lambda \subset \Conj $. Then $\dim \bigcap_{C \in \Lambda} \SC =1$ if and only if  $\bigcap_{C \in \Lambda} \SC=\Span \{I\}$.
\end{lemma}
\begin{proof}
We assume that $\dim \bigcap_{C \in \Lambda} \SC =1$. Then Lemma \ref{I} shows that $I \in \bigcap_{C \in \Lambda} \SC $. Since $\bigcap_{C \in \Lambda} \SC$ is a complex linear subspace of $B(\cH)$, we have $\Span \{I\} \subset \bigcap_{C \in \Lambda} \SC $. As $\dim \bigcap_{C \in \Lambda} \SC =1$, we obtain $\bigcap_{C \in \Lambda} \SC=\Span \{I\}$. The converse is clear.
\end{proof}

\begin{lemma}\label{projectionC}
Let $P \in \BH$ be a projection and $C \in \Conj$. Then 
\[
C(P(\cH))=P(\cH) \Longleftrightarrow P \in \SC.
\]
\end{lemma}
\begin{proof}
Assume that $C(P(\cH))=P(\cH)$. Let $x \in P(\cH)$ and  $y=C(x) \in P(\cH)$. Therefore, $CPCx=CPy=Cy=x=Px$. For any $x \in P(\cH)^{\perp}$, we have
\[
\langle Cx, y \rangle=\langle Cy, x \rangle =0
\]
for any $y \in P(\cH)$. This implies that $Cx \in  P(\cH)^{\perp}$. Hence we get $CPCx=0=Px$. Thus,  we have
\[
CPC=P.
\]
To prove the converse implication, let $x \in P(\cH)$. We have $Cx=CPx=PCx \in P(\cH)$. Thus $C(P(\cH)) \subset P(\cH)$. It follows that $P(\cH)=C^2(P(\cH)) \subset CP(\cH)$. Hence we get $C(P(\cH)) = P(\cH)$.
\end{proof}

\begin{remark}\label{projC}
 If $C(P(\cH)) = P(\cH)$, the restrictions of $C$ to $P(\cH)$ and  $(I-P)(\cH)$ are conjugations on $P(\cH)$ and $(I-P)(\cH)$, respectively.
\end{remark}

\begin{lemma}\label{onep}
Let $P \in \BH$ be a projection and $\Lambda=\{ C \in \Conj \mid C(P(\cH))=P(\cH)\}$. Then 
\[
\bigcap_{C \in \Lambda} \SC=\Span \{P, I\}.
\]
\end{lemma}
\begin{proof}
Applying Lemma \ref{projectionC}, we have $P \in \bigcap_{C \in \Lambda} \SC$.  It follows that 
\[
\Span\{P,I\} \subset \bigcap_{C \in \Lambda} \SC.
\]
To prove the reverse inclusion,  let $A \in \bigcap_{C \in \Lambda} \SC$. Since $\SC$ is closed under the involution $*$, we may assume that $A$ is a self-adjoint. 
Let $\cH_1=P(\cH)$ and $\cH_2=(I-P)(\cH)$. Let $\{e_j\}$ be an orthonormal basis of $\cH_1$ and $\{f_j\}$ be an orthonormal basis of $\cH_2$.  Fix $e_k \in \{e_j\}$. 
 Then there is $\alpha_j, \beta_j \in \mathbb{C}$ such that 
\begin{equation}\label{11}
Ae_k=\Sigma_{j} \alpha_j e_j+\Sigma_{j} \beta_j f_j. 
\end{equation}
Let $C$ be a conjugation associated with an orthonormal basis $\{ e_j, f_j \}$. Then $C \in \Lambda$. By \eqref{11}, we have 
\[
\Sigma_{j} \alpha_j e_j+\Sigma_{j} \beta_j f_j=Ae_k=CACe_k=CAe_k=\Sigma_{j} \overline{\alpha_j} e_j+\Sigma_{j} \overline{\beta_j }f_j. 
\]
Thus we have
\begin{equation}\label{112}
\alpha_j=\overline{\alpha_j}, \quad  \beta_j=\overline{\beta_j }, \quad \text{for all } j.
\end{equation}
Let $C'$ be a conjugation associated with a orthonormal basis $\{ (i)^{1-\delta_{(k,j)}}e_j, i f_j \}$. Then $C' \in \Lambda$. By \eqref{11}, we have 
\begin{equation*}
Ae_k=\alpha_k  e_k+\Sigma_{j\neq k} -i \alpha_j i e_j+\Sigma_{j} -i \beta_j if_j. 
\end{equation*}
Since $C'AC'=A$, we get 
\begin{multline*}
\Sigma_{j} \alpha_j e_j+\Sigma_{j} \beta_j f_j=Ae_k=C'Ae_k=C'(\alpha_k  e_k+\Sigma_{j\neq k} -i \alpha_j i e_j+\Sigma_{j} -i \beta_j if_j)\\=\overline{\alpha_k}e_k+\Sigma_{j\neq k} i \overline{\alpha_j }  i e_j +\Sigma_{j} i \overline{\beta_j } if_j=\overline{\alpha_k}e_k-\Sigma_{j\neq k} \overline{\alpha_j} e_j-\Sigma_{j} \overline{\beta_j }f_j.
\end{multline*}
This implies that %$\alpha_k=\overline{\alpha_k}$, $\alpha_j=-\overline{\alpha_j }$ for all  $j \neq k$ and $\beta_j=-\overline{\beta_j}$ for any $ j $. 
\begin{equation}\label{113}
\alpha_k=\overline{\alpha_k}, \quad  \alpha_j=-\overline{\alpha_j } \quad \text{for all} \ j \neq k , \quad \beta_j=-\overline{\beta_j} \  \text{for all } j.
\end{equation}
By \eqref{112} and \eqref{113}, we $\alpha_j=0$ for all  $j \neq k$ and $\beta_j=0$  for any $ j $. By \eqref{11}, we get
$Ae_k= \alpha_k e_k$.
This implies that 
\begin{equation}\label{ch1}
A(\cH_1) \subset \cH_1.
\end{equation}
Thus $A|_{\cH_1} \in  B(\cH_1)_{sa}$. By Remark \ref{projC}, $\{ C|_{\cH_1} \mid C \in \Lambda \}$ is the set of all conjugations on $\cH_1$. It implies that for any conjugation $C$ on $\cH_1$, we have $CA|_{\cH_1}C=A|_{\cH_1}$. Applying  Lemma \ref{I} , there is $\alpha \in \mathbb{R}$ such that 
\[
A|_{\cH_1} =\alpha I|_{\cH_1}.
\] 
By a similar argument for $\cH_2$,  we have 
\begin{equation}\label{ch2}
A(\cH_2) \subset \cH_2.
\end{equation}
Hence $A|_{\cH_2} \in  B(\cH_2)_{sa}$. 
Note that $\{ C|_{\cH_2} \mid C \in \Lambda \}$ is the set of all conjugations on $\cH_2$ by Remark \ref{projC}. Hence for any conjugation $C$ on $\cH_2$, we have $CA|_{\cH_2}C=A|_{\cH_2}$. Applying  Lemma \ref{I} , there is $\beta \in \mathbb{R}$ such that 
\[
A|_{\cH_2} =\beta I|_{\cH_2}.
\]
Therefore, we get 
\[
A=\alpha P+\beta (I-P) \in \Span \{P, I\}.
\] 
\end{proof}
By an argument similar to that used in the proof of Lemma \ref{onep}, we obtain the following.
\begin{lemma}\label{twop}
    Let $P_1, P_2 \in \BH$ be projections such that $P_1(\cH) \perp P_2(\cH)$ and $\Lambda=\{ C \in \Conj \mid CP_1(\cH)=P_1(\cH), CP_2(\cH)=P_2(\cH)\}$. Then 
    \[
     \bigcap_{C \in \Lambda} \SC =\Span\{ P_1, P_2, I\}.
    \]
\end{lemma}
\begin{definition}
  For any $x \in \cH \setminus \{0\} $, %with $\|x\|=1$, 
we denote 
\[
\Lambda_{x}=\{ C \in \Conj \mid C(x) \in \mathbb{T}x\}.
\]  
\end{definition}


In particular, we have the following.
\begin{cor}\label{dim2o}
Let $x \in \cH$ with $\|x\|=1$. Then 
\[
\bigcap_{C \in \Lambda_x} \SC=\Span\{x\otimes x, I\}.
\]
\end{cor}

\begin{lemma}\label{dim2t}
Let $\Lambda \subset \Conj$. Then  $\dim (\bigcap_{C \in \Lambda} \SC)=2$  if and only if 
there is a projection $P \in \BH$ with $P \neq I, 0$ such that %$C(P(\cH)) =P(\cH)$ for any $C \in \Lambda$ and 
$\bigcap_{C \in \Lambda} \SC=\Span \{P, I\}$.
\end{lemma}
\begin{proof}
Suppose that $\dim (\bigcap_{C \in \Lambda} \SC)=2$. By Lemma \ref{I}, $I \in \bigcap_{C \in \Lambda} \SC$. There is $A \in \bigcap_{C \in \Lambda} \SC$ such that $A$ is linearly independent of $I$. Since $\bigcap_{C \in \Lambda} \SC$ is closed under the involution $*$, we may assume that $A$ is a self-adjoint.
Suppose that $|\sigma(A)| \ge 3$. Let $\lambda_1, \lambda_2, \lambda_3 \in \sigma(A)$ be pairwise distinct. Since $\sigma(A)$ is a compact Hausdorff space, there are mutually disjoint open neighborhoods $U_i$ of $\lambda_i$ for $i=1,2,3$. By Urysohn's Lemma, for each $i=1,2,3$, there is a continuous function $f_i : \sigma(A) \to \mathbb{R}$ such that $f_i(\lambda_i)=1$,  $f_i(z)=0$ for all $z \in \sigma(A) \setminus U_i$ and $f_i(z) \ge 0$ for all $z \in \sigma(A)$. Note that $f_1f_2=f_2f_3=f_3f_1=0$. Applying the functional calculus we define a positive element $B_i=f_i(A) \in B(\cH)$ for each $i=1,2,3$. We get $B_1B_2=B_2B_3=B_3B_1=0$.  Suppose that 
\[
\alpha_1 B_1+ \alpha_2B_2+\alpha_3 B_3=0, \quad \text{for some} \ \alpha_1, \alpha_2, \alpha_3 \in \mathbb{C}.
\]
Multiplying both sides by $B_i$, we get 
\[
\alpha_iB_i^2=0.
\] 
By $\|B_i^2\|=\|B_i\|^2 \neq 0$, we obtain that $\alpha_i=0$. Hence $B_1$, $B_2$, and $B_3$ are linearly independent. On the other hand,  since $A \in \bigcap_{C \in \Lambda} \SC$,  we have $B_i \in \bigcap_{C \in \Lambda} \SC$.  This leads to a contradiction with $\dim (\bigcap_{C \in \Lambda} \SC)=2 $. Hence we have $|\sigma(A)| \le 2$. Therefore there are $\lambda_1, \lambda_2 \in \mathbb{R}$, a projection $P$ such that 
\begin{equation*}\label{SA}
A=\lambda_1 P+ \lambda_2(I-P)=(\lambda_1-\lambda_2)P+ \lambda_2I.
\end{equation*}
As $A$ is linearly independent of $I$, we get $P \neq I, 0$.  For any $C \in \Lambda$, we have
\[
\lambda_1 CPC+ \lambda_2(I-CPC)=CAC=A=\lambda_1 P+ \lambda_2(I-P).
\]
It is easy to check that $CPC$ is a projection. By the uniqueness of the spectrum decomposition of $A$, we have $CPC=P$ for any $C \in \Lambda$. Thus $P \in \bigcap_{C \in \Lambda} \SC$. %Applying Lemma \ref{projectionC}, we have $C(P(\cH)) = P(\cH)$ for all $C \in \Lambda$. %Thus $P \in \bigcap_{C \in \Lambda} \SC$. Since $CPC=P$, $CP(x)=PC(x) \in P(\cH)$ for any $x \in \cH$. We have $C(P(\cH)) \subset P(\cH)$ and $P(\cH)=C^2(P(\cH)) \subset CP(\cH)$. Hence we get $C(P(\cH)) = P(\cH)$, which means that every $C \in \Lambda$ is  a conjugation on $P(\cH)$.
Since $\dim(\bigcap_{C \in \Lambda} \SC)=2$, we have 
\[
\bigcap_{C \in \Lambda} \SC=\Span \{P, I\}.
\]
To prove the converse implication, suppose that $\bigcap_{C \in \Lambda} \SC=\Span \{P, I\}$ with $P \neq I, 0$. Then  it is clear that  $\dim (\bigcap_{C \in \Lambda} \SC)=2$.
\end{proof}





\begin{lemma}\label{I2}
Let $\cH$ be a separable Hilbert space. Let $T: \BH \to \BH$ be a continuous bijective linear map  and $\psi: \Conj \to \Conj$ be a bijection satisfying $T(\SC) = {\mathcal{S}_{\psi(C)}}$. There is $c \in \mathbb{C} \setminus \{0\} $ such that $T(I)=cI$.
\end{lemma}
\begin{proof}
%Let $\Lambda=\{D \in \Conj \mid D=\psi(C) , C \in \Conj \}$.  
Since $\psi$ is a bijection,  Lemma \ref{I} shows that
\[
T(I) \in T( \bigcap_{C \in \Conj} \SC)=  \bigcap_{C \in \Conj} \mathcal{S_{\psi(C)}}=\Span\{I\}.
\]
As $T$ is a bijective linear map,  there is $c \in \mathbb{C} \setminus \{0\} $ such that $T(I)=cI$. 
\end{proof}

\begin{lemma}\cite[Lemma 1]{AOS}\label{L1}
Let $\{e_n\}$ be an orthonormal basis of $\cH$. Then 
$D_{\{e_n\}}=\bigcap_{C \in \Lambda_{\{e_n\}}} \mathcal{S}_{C}$.
\end{lemma}

\begin{lemma}\label{I3}
Let $\cH$ be a separable Hilbert space. Let $T: \BH \to \BH$ be a continuous bijective linear map  and  $\varphi: \mathcal{E}(\cH) \to \mathcal{E}(\cH)$ be a bijection satisfying $T(\mathcal{D}_{\{e_n\}}) = \mathcal{D}_{\varphi(\{e_n\})}$. There is $c \in \mathbb{C} \setminus \{0\} $ such that $T(I)=cI$.
\end{lemma}
\begin{proof}
%Let $\Lambda=\{D \in \Conj \mid D=\psi(C) , C \in \Conj \}$.  
By Lemma \ref{L1}, we have $\bigcap_{\{e_n\} \in \mathcal{E}(\cH)} D_{\{e_n\}}= \bigcap_{C \in \Conj} \SC=\Span\{I\}$. 
Since $\varphi$ is a bijection,  we have
\begin{multline*}
   T(I) \in T(\bigcap_{\{e_n\} \in \mathcal{E}(\cH)} D_{\{e_n\}})=  \bigcap_{\{e_n\} \in \mathcal{E}(\cH)} T(D_{\{e_n\}})\\= \bigcap_{\{e_n\} \in \mathcal{E}(\cH)} D_{\varphi(\{e_n\})}=\Span\{I\}. 
\end{multline*}


As $T$ is a bijective linear map,  there is $c \in \mathbb{C} \setminus \{0\} $ such that $T(I)=cI$. 
\end{proof}

At the end of this section, we introduce an equivalence relation $\sim$ on $\mathcal{E}(\cH)$ by
\[
\{e_n\} \sim \{f_n\} 
\quad\Longleftrightarrow\quad
D_{\{e_n\}}=D_{\{f_n\}}
\]
for $\{e_n\},\{f_n\}\in \mathcal{E}(H)$. The equivalence class of $\{e_n\} \in \mathcal{E}(\cH)$ is denoted by $[e_n]$, and we write $\mathcal{E}(\cH)/ \sim$ for the corresponding quotient space. Let $\mathfrak{S}_N$ be the group of all permutations of  $\mathcal{I}_N$. We emphasize that, throughout this paper, an orthonormal basis is regarded as an ordered sequence. Thus, for an orthonormal basis $\{e_n\}$ and a permutation $\sigma \in \mathfrak{S}_N$, the sequences $\{e_{\sigma(n)}\}$ are regarded as distinct elements of $\mathcal{E}(\mathcal{H})$ unless $\sigma$ is the identity permutation. Hence, $\{e_n\} \sim \{f_n\} $ if and only if there exist a permutation $\sigma\in \mathfrak{S}_N$ and $\{\alpha_n\} \in \mathbb T^{N}$ such that $f_n=\alpha_n e_{\sigma(n)}$ for all $n$.  Let $q:\mathcal{E}(\cH) \to \mathcal{E}(\cH)/ \sim $ be the canonical quotient map. Fix a section $s:\mathcal{E}(\cH)/ \sim  \to \mathcal{E}(\cH)$ of $q$. 
\begin{lemma}\label{lift}
Suppose that $\mathcal{F}:\mathcal{E}(\cH)/ \sim \to \mathcal{E}(\cH)/ \sim$
is a bijection. Then there is a bijection $\varphi:\mathcal{E}(\cH) \to \mathcal{E}(\cH)$ such that $q\circ \varphi=\mathcal{F} \circ q$.
\end{lemma}
\begin{proof}
   %Let $G=\mathbb T^{\mathbb N}\rtimes S_\infty$. 
We define a left action of  the group $\mathbb T^{N}\rtimes \mathfrak{S}_N$ on $\mathcal{E}(\cH)$ by
\[
(\{\alpha_n\}, \sigma)   \cdot\{e_n\}=\{\alpha_n e_{\sigma(n)}\}, 
\]
for $(\{\alpha_n\}, \sigma) \in \mathbb T^{N}\rtimes \mathfrak{S}_N$ and $\{e_n\} \in \mathcal{E}(\cH)$. Then the group $\mathbb T^{N}\rtimes \mathfrak{S}_N$ acts freely and transitively on each equivalence class. Hence, every $\{e_n\} \in \mathcal{E}(\cH)$ can be written uniquely in the form
\[
\{e_n\}=g\cdot s(q(\{e_n\})).
\]
Consequently, we may identify $\mathcal{E}(\cH)$ with $(\mathcal{E}(\cH)/\sim) \times (\mathbb T^{N}\rtimes \mathfrak{S}_N)$.
Under this identification, define $\varphi:\mathcal{E}(\cH) \to  \mathcal{E}(\cH)$
by
\[
\varphi \bigl(g\cdot s(x)\bigr)=g\cdot s(F(x)), \qquad x \in\mathcal{E}(\cH)/\sim, \quad  g\in \mathbb T^{N}\rtimes \mathfrak{S}_N.
\]
Then $\varphi$ is a bijection and $q\circ \varphi=\mathcal{F} \circ q$.  
\end{proof}


\section{Proof of Theorem \ref{main0}} 
%The case where $\cH$ is two-dimensional requires a different argument, so we give it here first. 
Assume that $\dim \cH=2$ in this section.
We denote by
\[
U(2)=\{U\in M_2(\mathbb C):U^*U=I\}
\]
the unitary group of $2\times2$ complex matrices, and by
\[
SO(3)=\{R\in M_3(\mathbb R):R^TR=I,\ \det R=1\}
\]
the special orthogonal group of $\mathbb R^3$. %Thus, the elements of
%$U(2)$ preserve the Hermitian inner product on $\mathbb C^2$, while the
%elements of $SO(3)$ preserve the Euclidean inner product on $\mathbb R^3$
%and have determinant one.
\begin{proposition}{\cite[Proposition 5.5]{Land}}\label{Bloch_representation}
For each $U\in U(2)$, there exists a unique orthogonal transformation
$R_U\in SO(3)$ such that
\[
U(x\cdot \sigma)U^*=(R_Ux)\cdot\sigma
\qquad x \in\mathbb R^3,
\]
where
\[
x\cdot\sigma=\sum_{j=1}^3x_j\sigma_j.
\]
Moreover, the map $U \mapsto R_U$ is a group homomorphism from $U(2)$ onto $SO(3)$.  
\end{proposition}

\begin{lemma}\label{plane}
We have
\begin{multline*}
    \{\SC \mid C \in \Conj\}\\=\{\Span\{I, P\cdot\sigma\} \mid P \text{ is  a two-dimensional subspace of $\mathbb{R}^3$}\}
\end{multline*}
\end{lemma}
\begin{proof}
Let $K$ be a conjugation associated with the standard basis. For any $C \in \Conj$, there is $U \in U(2)$ such that $C=UKU^{*}$. (Let $\{f_i\}$ be an orthonormal basis such that $C(f_i)=f_i$. Define $U \in U(2)$ by $U(e_i)=f_i$ for $i=1,2$, where $\{e_i\}$ is the standard basis.) Then $\SC=U\mathcal{S}_KU^{*}$. We have $\mathcal{S}_K=\Span\{I, \sigma_1, \sigma_3\}$. Let $P_K=\{ (x_1,0,x_3) \in \mathbb{R}^3 \mid x_1, x_3 \in \mathbb{R}\}$ be a two-dimensional subspace of $\mathbb{R}^3$. It follows that 
\[
\mathcal{S}_K=\Span\{I, P_K \cdot \sigma\}.
\]
By Proposition \ref{Bloch_representation}, we have
\begin{equation*}
\begin{split}
\SC&=U\mathcal{S}_KU^{*}\\
&=U\Span\{I, P_K \cdot \sigma\}U^{*}=\Span\{I, U P_K \cdot \sigma U^{*}\}\\
&=\Span\{I, R_U P_K \cdot \sigma \}.
\end{split}
\end{equation*}
Since $R_U \in SO(3)$, $R_U P_K$ is  a two-dimensional subspace of $\mathbb{R}^3$. 


We shall the converse. Let $P$ be a two-dimensional subspace of $\mathbb{R}^3$. Then there is $z \in \mathbb{R}^3$ such that $P=\{x \in \mathbb{R}^3 \mid \la x, z\ra=0\}$. There is $R \in SO(3)$ such that $R(0,1,0)=z$. As  the elements of $SO(3)$ preserve the Euclidean inner product on $\mathbb R^3$, we have  
\[
RP_K=R(\{x \in \mathbb{R}^3 \mid \la x, (0,1,0)\ra =0\})=\{x \in \mathbb{R}^3 \mid \la x, z\ra=0\}=P.
\]
By Proposition \ref{Bloch_representation}, there is $U \in U(2)$ such that $R=R_U$. Define $D=UKU^{*}$. Then $D \in \Conj$ and 
\[
\mathcal{S}_D=U\mathcal{S}_KU^{*}=\Span\{I, R_U P_K \cdot \sigma \}=\Span\{I, P \cdot \sigma \}
\]
as the same argument above. 
\end{proof}

\begin{proof}[{\textbf {Proof of Theorem \ref{main0}}}] 
$\bm{(1) \Rightarrow (2)}.$ Let $\{e_1,e_2\} \in \mathcal{E}(\cH)$ and  $C \in \Lambda_{e_1}$. Corollary \ref{dim2o} shows that 
\[
\bigcap_{C \in \Lambda_{e_1}} \SC=\Span\{e_1 \otimes e_1, I \}=D_{\{e_1,e_2\}}. 
\]
Since $T(\bigcap_{C \in \Lambda_{e_1}} \SC)=\bigcap_{C \in \Lambda_{e_1}} \mathcal{S}_{\psi(C)}$ and $\dim \bigcap_{C \in \Lambda_{e_1}} \SC=2$, we have $\dim \bigcap_{C \in \Lambda_{e_1}} \mathcal{S}_{\psi(C)}=2$. By Lemma \ref{dim2t},  there is $e'_1 \in \cH$ with $\|e'_1\|=1$ such that 
\[
\bigcap_{C \in \Lambda_{e_1}} \mathcal{S}_{\psi(C)}=\Span\{e'_1 \otimes e'_1, I \}=D_{\{e'_1,e'_2\}},
\]
where $\{e'_1, e'_2\} \in \mathcal{E}(\cH)$. 
Hence, we have 
\begin{equation}\label{diago}
T(D_{\{e_1,e_2\}})=D_{\{e'_1,e'_2\}}, 
\end{equation}
for any $\{e_1,e_2\}  \in \mathcal{E}(\cH)$. Since a similar argument holds for $T^{-1}$, there is a bijective map $\mathcal{F}: \mathcal{E}(\cH)/ \sim \to \mathcal{E}(\cH)/ \sim$ such that 
\[
T(D_{[e_n]})=D_{\mathcal{F}([e_n])}.
\]
Applying Lemma \ref{lift}, there is a bijection $\varphi$ such that $q\circ \varphi=\mathcal{F} \circ q$. Therefore we have
\[
T(D_{\{e_n\}})=T(D_{[e_n]})=D_{\mathcal{F}([e_n])}=D_{[\varphi(\{e_n\})]}=D_{\varphi(\{e_n\})}.
\] 
This completes the proof of  $(1)\Rightarrow(2)$. \\
%As $x$ was chosen arbitrarily, \eqref{diago} 
%%%%%%%%%%%%%%%%%%%%%%%%%%%%%%%%%%%%%%
$\bm{(2) \Rightarrow (3)}.$ 
By Lemma \ref{I3}, there is $c \in \mathbb{C} \setminus \{0\}$ such that $T(I)=cI$. Since $T$ is a complex linear map, there are linear maps $\ell: \mathbb{C}^3 \to \mathbb{C}$ and $L: \mathbb{C}^3 \to \mathbb{C}^3$ such that 
\[
T(x \cdot \sigma)=\ell(x)I+L(x)\cdot \sigma.
\]
Since $T$ is a bijection, $L$ is an invertible matrix. For any $y=(y_1, y_2, y_3) \in \mathbb{R}^3 \setminus \{0\}$, there is $\{e_1, e_2\} \in \mathcal{E}(\cH)$ such that $\frac{1}{2}(I+\frac{y}{\|y\|}\cdot \sigma)=e_1\otimes e_1$. Then %For any rank one projection $e\otimes e$, there is a rank one projection $f \otimes f$ such that 
\[
T(\Span\{e_1\otimes e_1, I\})=\Span\{f_1 \otimes f_1, I\},
\]
where $\varphi(\{e_1, e_2\})=\{f_1, f_2\}$. Thus there is $z=(z_1, z_2, z_3) \in \mathbb{R}^3$ with $z_1^2+z_2^2+z_3^2=1$ such that $f_1 \otimes f_1=\frac{1}{2}(I+z \cdot \sigma)$. 
It implies 
\begin{multline*}
\Span\{I, L(y)\cdot \sigma\}=T(\Span\{I, y/\|y\| \cdot \sigma\})\\=T(\Span\{e_1\otimes e_1, I\})=\Span\{f_1 \otimes f_1, I\}=\Span\{I, z \cdot \sigma\}.
\end{multline*}
As $\{I, \sigma_1, \sigma_2, \sigma_3\}$ are linearly independent,  we have $L(y)=\lambda z$ for some $\lambda \in \mathbb{C}$. Let $e_1=(1,0,0)$ , $e_2=(0,1,0)$ and $e_3=(0,0,1)$ be the standard basis on $\mathbb{R}^3$. Then there are $b_i \in \mathbb{R}^3$ with $\|b_i\|=1$ such that $L(e_i)=\lambda_i b_i$ for some $\lambda_i \in \mathbb{C}$ for $i=1,2,3$. 
In addition as $e_1+e_2+e_3 \in \mathbb{R}^3$, there is $b \in \mathbb{R}^3$ such that $L(e_1+e_2+e_3)=\lambda b$ for some $\lambda \in \mathbb{C}\setminus\{0\}$. Since $L$  is an invertible matrix and $\{e_1, e_2, e_3\}$  is an orthonormal basis, there is $\{b_1, b_2, b_3\}$ is linearly independent. As $b_i \in \mathbb{R}^3$ and $b \in \mathbb{R}^3$, there are $r_i \in \mathbb{R}$ such that $b=r_1b_1+r_2b_2+r_3b_3$. It implies that
\[
\lambda_1 b_1+\lambda_2 b_2+\lambda_3 b_3=L(e_1+e_2+e_3)=\lambda b= \lambda (r_1b_1+r_2b_2+r_3b_3).
\]
We get $\lambda_i=\lambda r_i$ for any $i=1,2,3$. We define $B \in GL(3,\mathbb{R})$ by $Be_i=r_ib_i$ for any $i=1,2,3$. Then we have $L=\lambda B$. Therefore we have
\[
T(aI+x\cdot \sigma)=c aI+\ell(x)I+\lambda Bx.
\] 

$\bm{(3) \Rightarrow (1)}.$
Applying Lemma \ref{plane}, for any $C \in \Conj$, we have 
\[
\SC=\Span\{I, P\cdot\sigma\},  
\]
where $P$ is a  two-dimensional subspace of $\mathbb{R}^3$. For any $x \in P$, we have
\[
T(aI+x \cdot \sigma)=c aI+\ell(x)I+\lambda Bx \cdot \sigma \in \Span\{ I, BP \cdot \sigma\}.
\]
Since $B \in GL(3, \mathbb{R})$, we get $BP$  is a  two-dimensional subspace of $\mathbb{R}^3$. Hence there is $D \in \Conj$ such that $\Span\{ I, BP \cdot \sigma\}=\mathcal{S}_D$. As $T$ is a bijective linear map, we obtain $T(\SC)=\mathcal{S}_D$. The same argument applies to $T^{-1}$. Thus, by choosing a section of the quotient space of conjugations modulo the relation $C\sim e^{i\theta}C$, we can define, as above, a bijection $\psi:\Conj \to \Conj$ satisfying
\[
T(S_C)=S_{\psi(C)}.
\]

%%%%%%%%%%%%%%%%%%%%%%%%%%%%%%%%%%%%%%%%%%%

\end{proof}


\section{Proof of Theorem \ref{main}}
In this section, we assume that $\dim \cH \ge 3$. The implication $(3)\Rightarrow(1)$ is clear. Indeed, suppose that $(3)$ holds. Define $\psi(C)=UCU^{*}$ for any $C \in \Conj$. Then $\psi$ is a bijection on $\Conj$ that preserves commutativity. Moreover, we have  $T(\SC)=\mathcal{S}_{\psi(C)}$ for any $C \in \Conj$.
Thus, it remains to prove $(1)\Rightarrow(2)$ and $(2)\Rightarrow(3)$ in this paper.  

We first show that $(1)\Rightarrow(2)$.  For any orthonormal basis $\{e_i\}_{i \in \mathcal{I}_N}$ and a map $\rho: \mathcal{I}_N \to \{\pm1\} $, we define a conjugation $C_{\{e_i\}}^{\rho} \in \Conj$ by 
\[
C_{\{e_i\}}^{\rho}e_j= \rho(j) e_j, \quad \forall j \in \mathcal{I}_N.
\]
We define \[
\Lambda_{\{e_n\}}=\{C_{\{e_i\}}^{\rho} \in \Conj \mid \rho: \mathcal{I}_N \to \{\pm1\} \}.
\]
\begin{lemma}\label{commutative}
 Let $\psi: \Conj \to \Conj$ be a bijection that preserves commutativity. Then for any orthonormal basis $\{e_i\}_{i \in \mathcal{I}_N}$, there is an orthonormal basis $\{f_i\}_{i \in \mathcal{I}_N}$ such that 
 \[
 \psi(\Lambda_{\{e_n\}})=\Lambda_{\{f_n\}}.
 \]
\end{lemma}
\begin{proof}
   Let $\{e_i\}$ be an orthonormal basis.  Fix $C_0 \in \Lambda_{\{e_n\}}$. For any $C_{\{e_i\}}^{\rho} \in \Lambda_{\{e_n\}}$, we have $C_0C_{\{e_i\}}^{\rho}=C_{\{e_i\}}^{\rho}C_0$. Since $\psi$ preserves commutativity,  we have $\psi(C_0)\psi(C_{\{e_i\}}^{\rho})=\psi(C_{\{e_i\}}^{\rho})\psi(C_0)$. Define $U^{\rho}=\psi(C_0)\psi(C_{\{e_i\}}^{\rho})$. Then we have $(U^{\rho})^*=U^{\rho}$ and $(U^{\rho})^2=I$, thus $U^{\rho}$ is a self-adjoint operator with  $\sigma(U^{\rho}) \subset \{\pm 1\}$.  Let $\mathcal{U}=\{U^{\rho} \mid \rho: \mathcal{I}_N \to \{\pm1\} \}$. For any $U^{\rho_1}, U^{\rho_2} \in \mathcal{U}$, we get
   \[
   \begin{split}
        U^{\rho_1}U^{\rho_2}&=\psi(C_0)\psi(C_{\{e_i\}}^{\rho_1})\psi(C_0)\psi(C_{\{e_i\}}^{\rho_2})\\&=\psi(C_0)\psi(C_{\{e_i\}}^{\rho_2})\psi(C_0)\psi(C_{\{e_i\}}^{\rho_1})=U^{\rho_2}U^{\rho_1}.
   \end{split}
   \]
Hence $\mathcal{U}$ is commutative under multiplication. It implies that there is an orthonormal basis such that the elements of $\mathcal{U}$ are simultaneously diagonalizable with respect to the basis. For any map $\epsilon: \{ \rho: \mathcal{I}_N \to \{\pm1\}\} \to  \{\pm1\}$, we define $\cH_{\epsilon}=\bigcap_{\rho}\ker(U^{\rho}-\epsilon(\rho)I)$. Then we get $\cH=\oplus_{\epsilon} \cH_{\epsilon}$. We remark that
\begin{equation*}%\label{dec}
    x \in \cH_{\epsilon} \Longleftrightarrow U^{\rho}x=\epsilon(\rho)x, \quad \text{for any}\ \rho.
\end{equation*}
Let $x \in \cH_{\epsilon}$. Then we get
\[
U^{\rho}(\psi(C_0)x)=\psi(C_0)(U^{\rho}x)=\psi(C_0)\epsilon(\rho)x=\epsilon(\rho)\psi(C_0)x
\]
for any $\rho$. Thus we get $\psi(C_0)x \in \cH_{\epsilon}$. It follows that for any $\epsilon$,
\[
\psi(C_0)(\cH_{\epsilon})=\cH_{\epsilon}.
\]
Thus $\psi(C_0)|_{\cH_{\epsilon}}$ is a conjugation on $\cH_{\epsilon}$. Applying Lemma \ref{GP1}, there is an orthonormal basis $F_{\epsilon}$ on $\cH_{\epsilon}$ such that $\psi(C_0)|_{\cH_{\epsilon}}$ is a conjugation associated with $F_{\epsilon}$. Choose $\{f_n\} \in \mathcal{E}(\cH)$ such that $\bigcup_{\epsilon}F_{\epsilon} =\{f_n \mid n \in I_N\} $. It follows that $\psi(C_0)f_n=f_n$ for any $n \in I_N$. For any $f_n \in \cH_{\epsilon}$, we have
\[
\psi(C_{\{e_i\}}^{\rho})f_n=\psi(C_{\{e_i\}}^{\rho})\psi(C_0)f_n=U^{\rho}f_n=\epsilon(\rho)f_n.
\]
This implies that for any $\rho$, $\psi(C_{\{e_i\}}^{\rho}) \in \Lambda_{\{f_n\}}$. Therefore we have
\[
\psi(\Lambda_{\{e_n\}}) \subset \Lambda_{\{f_n\}}.
\]
Let $D \in \Lambda_{\{f_n\}}$. Since $\psi$ is bijective, there is $C \in \Conj$ such that $\psi(C)=D$.  As the elements of $\Lambda_{\{f_n\}}$ are commutative under multiplication, for any $\rho$ we have
\[
\psi(C_{\{e_i\}}^{\rho})\psi(C)=\psi(C)\psi(C_{\{e_i\}}^{\rho}).
\]
Since $\psi$ preserves commutativity, we get $C_{\{e_i\}}^{\rho}C=CC_{\{e_i\}}^{\rho}$.
Fix $n \in \mathcal{I}_N$. Then we have 
\[
C_{\{e_i\}}^{\rho}C(e_n)=CC_{\{e_i\}}^{\rho}(e_n)=C\rho(n)e_n=\rho(n)Ce_n.
\]
Take $\rho \equiv1$. Then we have $C_{\{e_i\}}^{1}Ce_n=Ce_n$. Thus there is $\alpha_i \in \mathbb{R}$ such that
\[
Ce_n=\sum_{i}\alpha_i e_i
\]
as $\{e_i\}$ is an orthonormal basis. On the other hand, when $\rho_n(i)=(-1)^{\delta_{n,i}}$,  we have $C_{\{e_i\}}^{\rho_n}C(e_n)=-C(e_n)$. Since
\[
C_{\{e_i\}}^{\rho_n}C(e_n)=C_{\{e_i\}}^{\rho_n}(\sum_{i}\alpha_i e_i)=\sum_{i}\alpha_iC_{\{e_i\}}^{\rho_n} e_i=\sum_{i}\alpha_i\rho_n(i) e_i=-\alpha_n e_n+\sum_{i \neq n} \alpha_i e_i
\]
and
\[
-C(e_n)=\sum_{i}-\alpha_i e_i,
\]
we get  $\alpha_{i}=0$  if $\ i \neq n$. Therefore 
\[
Ce_n=\alpha_ne_n.
\]
Since $C$ is an isometry, $\alpha_n \in \{ \pm1\}$.
As this holds for every $n \in \mathcal{I}_N$, we see that $C \in \Lambda_{\{e_n\}}$. This implies that $D=\psi(C) \in \psi(\Lambda_{\{e_n\}})$, and 
\[
\psi(\Lambda_{\{e_n\}})=\Lambda_{\{f_n\}}.
\]
\end{proof}



\begin{proof}[{\textbf {Proof of Theorem \ref{main}: $\bm{(1) \Rightarrow (2)}$}}] 
Let $\{e_n\}$ be an orthonormal basis of $\cH$. Applying Lemma \ref{commutative}, there is an orthonormal basis $\{f_n\}$ such that 
 \[
 \psi(\Lambda_{\{e_n\}})=\Lambda_{\{f_n\}}.
 \]
By Lemma \ref{L1},  we have 
 \begin{equation*}
 \begin{split}
      T(D_{\{e_n\}})&=T(\bigcap_{C \in \Lambda_{\{e_n\}}} \mathcal{S}_{C})=\bigcap_{C \in \Lambda_{\{e_n\}}} T(S_c)\\&=\bigcap_{C \in \Lambda_{\{e_n\}}} S_{\psi(C)}=\bigcap_{C \in \Lambda_{\{f_n\}}} S_C=D_{\{f_n\}}.
 \end{split}
 \end{equation*}
Since $T$ is a bijection and a similar argument as above holds for $T^{-1}$, there is a bijective map $\mathcal{F}:\mathcal{E}(\cH)/ \sim  \to \mathcal{E}(\cH)/ \sim $ by 
\[
T(D_{[e_n]})=D_{\mathcal{F}([e_n])}.
\]
Applying Lemma \ref{lift}, there is  a bijection $\varphi$ satisfying $q\circ \varphi=\mathcal{F} \circ q$. Therefore we have
\[
T(D_{\{e_n\}})=T(D_{[e_n]})=D_{\mathcal{F}([e_n])}=D_{[\varphi\{e_n\}]}=D_{\varphi(\{e_n\})}.
\] 
\end{proof}

%\subsection{Proofs for (2) $\Rightarrow$ (3)}
Now we show $(2)\Rightarrow(3)$. We assume that $T: \BH \to \BH$ is a continuous bijective linear map  and $\varphi: \mathcal{E}(\cH) \to \mathcal{E}(\cH)$ is a bijection satisfying 
\[
T(\mathcal{D}_{\{e_n\}}) = \mathcal{D}_{\varphi(\{e_n\})},
\] 
for any $\{e_n\} \in \mathcal{E}(\cH)$. 

\begin{lemma}\label{projection1}
Let $P_1, P_2 \in \BH $ be projections such that $P_1(\cH) \perp P_2(\cH)$. 
Then there are projections $Q_1, Q_2 \in \BH$
such that 
\[
T(\Span \{P_1, I\} )=\Span \{Q_1,I\},
\]
 and 
\[
T(\Span \{P_2, I\} )=\Span \{Q_2,I\}.
\]
Then, one of the following holds;
%\begin{enumerate}
%\item $Q_1Q_2=0$,
%\item $Q_1Q_2=Q_1$,
%\item $Q_1Q_2=Q_2$,
%\item $(I-Q_1)(I-Q_2)=0$.
%\end{enumerate}
\[
Q_1Q_2=0, \ Q_1Q_2=Q_1, \ Q_1Q_2=Q_2, \ \text{or} \ (I-Q_1)(I-Q_2)=0.
\]
\end{lemma}

\begin{proof}
Let $P_1, P_2 \in \BH$ be projections such that $P_1(\cH) \perp P_2(\cH)$. 
For any $i=1,2$, we define $\mathcal{E}_i=\{ \{e_n\} \in \mathcal{E}(\cH) \mid \{e_n\} \cap P_i(\cH) \in \mathcal{E}(P_i(\cH))\}$. Then we shall show that 
\begin{equation}\label{cap1}
    \{ C \in \Conj \mid C \in \Lambda_{\{e_n\}}, \{e_n\} \in \mathcal{E}_i\}=\{ C \in \Conj \mid CP_i(\cH)=P_i(\cH)\}.
\end{equation}
Let $C \in \Lambda_{\{e_n\}}$ with $\{e_n\} \in \mathcal{E}_i$. For any $x \in P_i(\cH)$, there is $\alpha_n \in \mathbb{C}$ such that $x=\sum_{\{n \mid e_n \in P_i(\cH)\}} \alpha_ne_n$. We have $C(x)=C(\sum_{\{n \mid e_n \in P_i(\cH)\}} \alpha_ne_n)=\sum_{\{n \mid e_n \in P_i(\cH)\}}\overline{\alpha_n}e_n \in P_i(\cH)$. Thus $CP_i(\cH)=P_i(\cH)$. \\
To prove the opposite inclusion, let $C \in \Conj$ with $CP_i(\cH)=P_i(\cH)$. It follows that $C|_{P_i(\cH)}$ is a conjugation on $P_i(\cH)$. Therefore there is $\{e_n\} \in \mathcal{E}_i$ such that $C(e_n)=e_n$. Then $C \in \Lambda_{\{e_n\}}$. 
By a similar argument, we obtain
\begin{equation}\label{cap12}
    \begin{split}
        &\{ C \in \Conj \mid C \in \Lambda_{\{e_n\}}, \{e_n\} \in \mathcal{E}_1 \cap \mathcal{E}_2\}\\&=\{ C \in \Conj \mid CP_1(\cH)=P_1(\cH), CP_2(\cH)=P_2(\cH)\}.
    \end{split}
\end{equation}
By Lemma \ref{onep}, \ref{L1} and \eqref{cap1}, we have 
\begin{equation*}
    \begin{split}
        \Span\{P_i, I\}&=\bigcap_{\{C \mid CP_i(\cH)=P_i(\cH)\}} \SC\\&
        =\bigcap_{\{e_n\} \in \mathcal{E}_i} \bigcap_{C \in \Lambda_{\{e_n\}}} \SC
        =\bigcap_{\{e_n\} \in \mathcal{E}_i} D_{\{e_n\}}. 
    \end{split}
\end{equation*}
Thus we get
\begin{equation*}
    \begin{split}
        T(\Span\{P_i, I\})
        &=T(\bigcap_{\{e_n\} \in \mathcal{E}_i} D_{\{e_n\}})\\&=\bigcap_{\{e_n\} \in \mathcal{E}_i} D_{\varphi(\{e_n\})}=\bigcap_{\{e_n\} \in \mathcal{E}_i} \bigcap_{C \in \Lambda_{\varphi(\{e_n\})}}\SC. 
    \end{split}
\end{equation*}
Since $\dim (\bigcap_{\{e_n\} \in \mathcal{E}_i} \bigcap_{C \in \Lambda_{\varphi(\{e_n\})}}\SC) \le 2$,  Lemma \ref{dim1} and Lemma \ref{dim2t} show that there are projections $Q_i \in \BH$ such that 
\[
T(\Span\{P_i, I\})=\Span\{Q_i, I\}, \quad i=1,2.
\]
Note 
\begin{equation}\label{Q1Q2}
    Q_i \in \bigcap_{\{e_n\} \in \mathcal{E}_i} \bigcap_{C \in \Lambda_{\varphi(\{e_n\})}}\SC \subset \bigcap_{\{e_n\} \in \mathcal{E}_1 \cap \mathcal{E}_2}\bigcap_{C \in \Lambda_{\varphi(\{e_n\})}}\SC.
\end{equation}
On the other hand, since $P_1(\cH) \perp P_2(\cH)$, $\mathcal{E}_1 \cap \mathcal{E}_2 \neq \emptyset$. We have

\begin{equation*}
    \begin{split}
&\bigcap_{\{e_n\} \in \mathcal{E}_1 \cap \mathcal{E}_2}\bigcap_{C \in \Lambda_{\varphi(\{e_n\})}}\SC\\&=\bigcap_{\{e_n\} \in \mathcal{E}_1 \cap \mathcal{E}_2}D_{\varphi(\{e_n\})}=\bigcap_{\{e_n\} \in \mathcal{E}_1 \cap \mathcal{E}_2}T(D_{\{e_n\}})\\&=T(\bigcap_{\{e_n\} \in \mathcal{E}_1 \cap \mathcal{E}_2}D_{\{e_n\}})=T(\bigcap_{\{e_n\} \in \mathcal{E}_1 \cap \mathcal{E}_2}\bigcap_{C \in \Lambda_{\{e_n\}}}\SC)\\&=T(\bigcap \{\SC \mid  C \in \Conj, CP_1(\cH)=P_1(\cH), CP_2(\cH)=P_2(\cH)\})\\&=T(\Span\{ P_1, P_2, I\})=\Span\{ Q_1, Q_2, I\}
        \end{split}
\end{equation*}
%\end{comment}
\begin{comment}
\begin{multline*}
\bigcap_{\{e_n\} \in \mathcal{E}_1 \cap \mathcal{E}_2}\bigcap_{C \in \Lambda_{\varphi(\{e_n\})}}\SC=\bigcap_{\{e_n\} \in \mathcal{E}_1 \cap \mathcal{E}_2}D_{\varphi(\{e_n\})}=\bigcap_{\{e_n\} \in \mathcal{E}_1 \cap \mathcal{E}_2}T(D_{\{e_n\}})\\=T(\bigcap_{\{e_n\} \in \mathcal{E}_1 \cap \mathcal{E}_2}D_{\{e_n\}})=T(\bigcap_{\{e_n\} \in \mathcal{E}_1 \cap \mathcal{E}_2}\bigcap_{C \in \Lambda_{\{e_n\}}}\SC)\\=T(\bigcap \{\SC \mid  C \in \Conj, CP_1(\cH)=P_1(\cH), CP_2(\cH)=P_2(\cH)\})\\=T(\Span\{ P_1, P_2, I\})=\Span\{ Q_1, Q_2, I\}
\end{multline*}
\end{comment}
by \eqref{cap12} and Lemma \ref{twop}. 
Let $\{e_n\} \in \mathcal{E}_1 \cap \mathcal{E}_2$. Then we get $\Span\{ P_1, P_2, I\} \subset D_{\{e_n\}}$. Therefore 
\[
\Span\{ Q_1, Q_2, I\} =T(\Span\{ P_1, P_2, I\}) \subset T(D_{\{e_n\}})=D_{\varphi(\{e_n\})}.
\]
This implies that $Q_1Q_2=Q_2Q_1$. 
Fix $C \in \Lambda_{\varphi(\{e_n\})}$ where $\{e_n\} \in \mathcal{E}_1 \cap \mathcal{E}_2$. By \eqref{Q1Q2}, we get $CQ_iC=Q_i$ for $i=1,2$. Hence we get
\[
CQ_1Q_2C=CQ_1CCQ_2C=Q_1Q_2=Q_2Q_1=(Q_1Q_2)^*.
\]
It follows that 
\[
Q_1Q_2 \in \bigcap_{\{e_n\} \in \mathcal{E}_1 \cap \mathcal{E}_2}\bigcap_{C \in \Lambda_{\varphi(\{e_n\})}}\SC=\Span\{ Q_1, Q_2, I\}.
\]
Suppose that $Q_1Q_2 \neq 0$. There are $\alpha, \beta, \gamma \in \mathbb{C}$ such that
\[
Q_1Q_2=\alpha Q_1 +\beta Q_2+\gamma I.
\]
Multiplying from the left by $Q_1$, we obtain
%\[
%Q_1Q_2=\alpha Q_1 +\beta Q_1Q_2+\gamma Q_1 .
%\]
%and 
\[
(1-\beta)Q_1Q_2=(\alpha+\gamma)Q_1.
\]
Multiplying from the right by $Q_2$, we get $(1-\beta)Q_1Q_2=(\alpha+\gamma)Q_1Q_2$. Since $Q_1Q_2 \neq 0$, $1-\beta=\alpha+\gamma$. We consider two cases; (i)$\beta \neq 1$ and (ii) $\beta =1$.\\
\begin{itemize}
\item[(i)]When $\beta \neq 1$, we get 
\[
Q_1Q_2=Q_1.
\]

\item[(ii)] When $\beta =1$, we have $\gamma=-\alpha$. Thus 
\[
Q_1Q_2=\alpha Q_1 + Q_2-\alpha I.
\]
\begin{itemize}
\item[(1)] When $\alpha \neq 1$, multiplying from the right by $Q_2$, we obtain $Q_1Q_2=\alpha Q_1Q_2 + Q_2-\alpha Q_2 $. It follows that
\[
(1-\alpha) Q_1Q_2=(1-\alpha) Q_2. 
\]
Since $\alpha \neq 1$, we get
\[
Q_1Q_2=Q_2.
\]
\item[(2)]When $\alpha =1$, we have $Q_1Q_2=Q_1 + Q_2- I $, which implies that
\[
(I-Q_1)(I-Q_2)=0.
\]
\end{itemize}
\end{itemize}
\end{proof}

\begin{comment}
On the other hand, if $Q_1Q_2=Q_1$, then $Q_2(\cH) \subset Q_1(\cH)$. Since $Q_1, Q_2, I \in \bigcap_{C \in \Lambda_0} \mathcal{S}_{\psi(C)}$  are linearly independent, $Q_2(\cH) \neq Q_1(\cH)$. Let $Q_3(\cH)$ be the orthogonal complement of $Q_2(\cH)$ in $Q_1(\cH)$ and $Q_3 \in \BH$ be the projection on $Q_3(\cH)$. Then $Q_1(\cH)=Q_2(\cH) \oplus Q_3(\cH)$ and $Q_3 \neq 0, I$. Since $Q_1, Q_2 \in \bigcap_{C \in \Lambda_0} \mathcal{S}_{\psi(C)}$, $\psi(C)$ is a conjugation on $Q_1(H)$ and $Q_2(H)$. Therefore $\psi(C)$ is also a conjugation on $Q_3(\cH)$, which implies that $Q_3 \in \bigcap_{C \in \Lambda_0} \mathcal{S}_{\psi(C)}$. Suppose that 
\[
\alpha Q_1+\beta Q_2+\gamma Q_3+\delta I=0, \quad \text{for some} \alpha, \beta, \gamma, \delta \in \mathbb{C}.
\]
For any $x \in Q_1(\cH)^{\perp} \setminus \{0\}$, we have $0=\delta x$. Thus $\delta=0$. Multiplying from the right by $Q_2$, we have $\alpha Q_1+\beta Q_2=0$. For any  $x \in Q_3(\cH)^{\perp} \setminus \{0\}$, we get $\alpha x=0$, which implies $\alpha=0$. For any $x \in Q_3(\cH) \setminus \{0\}$, we get $\gamma x=0$, which implies $\gamma=0$. Therefore $\beta=0$. Hence $Q_1, Q_2, Q_3, I$ are linearly independent, which is a contradiction to $\dim \bigcap_{C \in \Lambda_0} \mathcal{S}_{\psi(C)}=3$. Therefore $Q_1Q_2=Q_1$ does not occur. As the same arguement, we conclude that $Q_1Q_2=Q_2$ also does not occur. Therefore we get 
\[
Q_1Q_2=0 \quad \text{or} \quad (I-Q_1)(I-Q_2)=0.
\]
\end{comment}

We denote the set of all rank one projections on $\cH$ by $\PH$. 
\begin{lemma}\label{comm}
Let $P_1, P_2 \in \PH$ with $P_1P_2=0$. Then there are two projections $Q_1, Q_2 \in \PH$ %which satisfies either $Q_1$ or $I-Q_1$ if of rank $1$ and either $Q_2$ or $I-Q_2$ is of rank $1$ 
such that $Q_1Q_2=0$ and $T(P_1) \in \operatorname{span}\{Q_1, I\}$, $T(P_2) \in \operatorname{span}\{Q_2, I\}$. 
\end{lemma}
\begin{proof}
Since $P_1P_2=0$, by Lemma \ref{projection1}, there are projections $Q_1, Q_2 \in \BH \setminus \{I,0\}$ such that $Q_1Q_2=0$. Fix $i=1,2$. We denote $\mathcal{C}(P_i)=\{ A \in \BH \mid P_iA=AP_i\}$, which is the commutator for  $P_i$. We shall show that $T(\mathcal{C}(P_i)) \subset \mathcal{C}(Q_i)$. Let $B \in \mathcal{C}(P_i)$. We assume that $B$ is a self-adjoint operator. Then there is a spectrum measure $P$ on $\sigma(B)$ such that $B=\int \lambda dP$. Since $BP_i=P_iB$,  for any Borel set $E \subset \sigma(B)$,  $P_i$ commutes with $P(E)$. Hence $P_i P(E)$ is a projection for any $E$.   Fix $E \subset \sigma(B)$. Assume that $P_i P(E) \neq 0$. As $P_i$ is  of rank $1$, we have $P_i(\cH)  \subset P(E)(\cH)$. In this case, $P_i (\cH) \perp (I-P(E))(\cH)$. Otherwise we have $P_i P(E) =0$, which means $P_i(\cH) \perp P(E)(\cH)$. In either case,  by Lemma \ref{projection1}, there is a projection $Q(E) \in \BH $
such that  
\[
T(\Span\{P(E) , I\})=\Span \{ Q(E), I\},
\]
which satisfies $Q_iQ(E)=0$ or $Q_iQ(E)=Q_i$ or $Q_iQ(E)=Q(E)$ or $(I-Q_i)(I-Q(E))=0$. Therefore  we get 
\[
Q(E) \in \mathcal{C}(Q_i), \quad \text{for all} \ E. 
\]
We define the inclusion map $\operatorname{id}: \sigma(B) \to \mathbb{C}$ by $\operatorname{id}(\lambda)=\lambda$ for any $\lambda \in \sigma(B)$. 
Let $f_n=\sum_{i=1}^{k_n} a_{n_i} \chi_{E_{n_i}}$ be a simple function on $\sigma(B)$ which satisfies $f_n (\lambda) \le \operatorname{id}(\lambda) $ for any $\lambda \in \sigma(B)$ and $\|f_n-\operatorname{id}\|_{\infty} \to 0$ in $B_{\infty}(\sigma(B))$, which is the Banach algebra of all bounded complex-valued measurable functions on $\sigma(B)$. Let $B_n=\int f_n dP=\sum_{i=1}^{k_n} a_{n_i} P(E_{n_i})$. Then we have
\[
T(B_n)=T(\sum_{i=1}^{k_n} a_{n_i} P(E_{n_i}))=\sum_{i=1}^{k_n} a_{n_i} T(P(E_{n_i}))=\sum_{i=1}^{k_n} a_{n_i} (\alpha_{n_i}Q(E_{n_i})+\beta_{n_i} I)
\]
for some $\alpha_{n_i}, \beta_{n_i} \in \mathbb{C}$. This implies that 
\[
T(B_n) \in \mathcal{C}(Q_i)
\]
for all $n$. Since $B_n \to B$ in $\BH$ and $T$ is a continuous map, we get $T(B_n) \to T(B)$. As $\mathcal{C}(Q_i)$ is closed subspace of $\BH$, we get $T(B) \in \mathcal{C}(Q_i)$. Hence we get 
\[
T(\{ B \in \mathcal{C}(P_i) \mid B=B^{*}\} ) \subset \mathcal{C}(Q_i).
\]
For any $B \in \mathcal{C}(P_i)$. Note that $B_1:=(B+B^{*})/2 \in \mathcal{C}(P_i)$ and $B_2:=(B-B^{*})/2i \in \mathcal{C}(P_i)$ since $P_i=P_i^{*}$. Since  $T(B)Q_i=T(B_1+iB_2)Q_i=T(B_1)Q_i+iT(B_2)Q_i=Q_iT(B_1)+Q_iiT(B_2)=Q_iT(B)$, we get
\[
T( \mathcal{C}(P_i)  ) \subset \mathcal{C}(Q_i).
\]
Applying \cite[Lemma 3.2]{Oml},  the subspace $ \mathcal{C}(P_1)+ \mathcal{C}(P_2)$ of $\BH$ has codimension $2$ in $\BH$. Since we have
\[
T(\mathcal{C}(P_1)+ \mathcal{C}(P_2)) \subset \mathcal{C}(Q_1) + \mathcal{C}(Q_2), 
\]
the codimension of the subspace $ \mathcal{C}(Q_1)+ \mathcal{C}(Q_2)$ of $\BH$ is less than or equal to $2$.  
 Suppose that the rank of $Q_1$ is greater than $1$.  Choose linearly independent vectors $x_1, x_2 \in Q_1(\cH)$ with $x_1 \perp x_2$ and $x_3 \in Q_2(\cH)$. Then the image of the quotient map of the operators
\[
x_1 \otimes x_3, \quad x_2 \otimes x_3, \quad x_3 \otimes x_1, \quad x_3 \otimes x_2
\]
are linearly independent operators of $\BH/\mathcal{C}(Q_1) + \mathcal{C}(Q_2)$. This is a contradiction to codimension of $ \mathcal{C}(Q_1)+ \mathcal{C}(Q_2)$ of $\BH$ is less than or equal to $2$. Thus the rank of $Q_1$ is $1$ and similarly the rank of $Q_2$ is $1$.
\end{proof}

\begin{theorem}[Uhlhorn's theorem \cite{UHL}]
    Suppose that $\dim \cH \ge 3$. Let $\Phi: \PH \to \PH$ be a bijective map such that 
    \[
P_1P_2=0 \Longleftrightarrow \Phi(P_1)\Phi(P_2)=0
    \]
    for any $P_1, P_2 \in \PH$. Then there is a unitary operator or antiunitary operator $U \in \BH$ such that $\Phi(P)=UPU^{*}$ for all $P \in \PH$. 
    
\end{theorem}

\begin{lemma}\label{UHL}
There is a unitary operator or antiunitary operator $U$ such that 
\[
T(P) \in \operatorname{span}\{UPU^{*}, I\} , \quad P \in \PH.
\]
\end{lemma}

\begin{proof}
Let $P \in \PH$.  Lemma \ref{comm} shows there is $Q_1 \in \PH$ such that $T(P) \in \operatorname{span}\{Q_1, I\}$. To show  the uniqueness of  $Q_1$,  suppose that there are $Q_1, Q_2 \in \PH$ such that 
$T(P)\in \operatorname{span}\{Q_1, I\}=\operatorname{span}\{Q_2, I\}$. Then there are $\alpha_i, \beta_i \in \mathbb{C}$ such that
\[
\alpha_1 Q_1+\beta_1I=\alpha_2 Q_2+\beta_2I=T(P).
\]
Since $\dim \cH \ge 3$, there is non-zero vector $y \in \operatorname{span}\{ Q_1(\cH), Q_2(\cH)\}^{\perp}$. We have $\beta_1y=\beta_2y$. Thus $\beta_1=\beta_2$. As $Q_1$ and $Q_2$ are non-trivial projections, we get $Q_1=Q_2$. Therefore we can define  $\Phi: \PH \to \PH$ by $T(P) \in \operatorname{span}\{\Phi(P) ,I \} $ for any $P \in \PH$. Let $P_1,P_2 \in \PH $ with $P_1P_2=0$. Applying Lemma \ref{comm},  we get
\[
\Phi(P_1)\Phi(P_2)=0.
\] 
By a similar argument as above for $T^{-1}$,  we also obtain 
\[
\Phi(P_1) \Phi(P_2)=0  \Longrightarrow P_1P_2=0.
\]
Applying Uhlhorn's theorem, there is a unitary operator or antiunitary operator $U \in \BH$ such that $\Phi(P)=UPU^{*}$ for all $P \in \PH$. 
\end{proof}

\begin{proof}[{\textbf {Proof of Theorem \ref{main}: $\bm{(2) \Rightarrow (3)}$}}] 
Let $\{e_n\} \in \mathcal{E}(\cH)$. Then  $e_i \otimes e_i \in \mathcal{D}_{\{e_n\}}$. Thus we have $T(e_i \otimes e_i) \in \mathcal{D}_{\varphi(\{e_n\})}$. By Lemma \ref{UHL} we have $T(e_i \otimes e_i) \in \Span \{U(e_i \otimes e_i)U^{*}, I\} =\Span \{U(e_i) \otimes U(e_i), I\} $. Thus $U(e_i) \otimes U(e_i) \in \mathcal{D}_{\varphi(\{e_n\})}$. This implies that 
$D_{\{U(e_i)\}} \subset \mathcal{D}_{\varphi(\{e_n\})}$.
On the other hand,  $\{U(e_i)\}$ and $\varphi(\{e_n\})$ are orthonormal bases of $\cH$. Hence
\[
D_{\{U(e_i)\}} = \mathcal{D}_{\varphi(\{e_n\})}.
\]
Define a bijective linear map $\widetilde{T}$ on $B(\cH)$ by $A \mapsto U^{*}T(A)U$. Then 
\[
\widetilde{T}(D_{\{e_n\}})=U^{*}T(D_{\{e_n\}})U=U^{*}D_{\{U(e_i)\}}U=D_{\{e_n\}}
\]
holds for any $\{e_n\} \in \mathcal{E}(\cH)$. By applying Theorem \ref{AOST}, there exists $c \in \mathbb{C} \setminus \{0\}$, a bounded linear functional $f$ on $\BH$ such that 
\[
\widetilde{T}(A)=cA+f(A)I, \quad A \in \BH.
\]   
Therefore we have
\[
T(A)=cUAU^{*}+f(A)I, \quad A \in \BH.
\]
\end{proof}



\begin{thebibliography}{99}
\bibitem{AOS} Z. Amara, M. Oudghiri and K. Souilah, Complex symmetric operators and additive preservers problem, Adv. Oper. Theory {\bf 5} (2020), no.~1, 261--279.

\bibitem{Brs} M. Bre\v{s}ar, Normal preserving linear mappings. Can. Math. Bull. 37, 306–309 (1994)

%\bibitem{Conway} J. B. Conway,  A course in functional analysis, Springer, New York (1997).


\bibitem{GP} S. R. Garcia and M. Putinar,  Complex symmetric operators and applications, Trans. Am. Math. Soc. 358, No. 3, 1285--1315.

\bibitem{Land} K. Landsman,  Foundations of quantum theory. From classical concepts to operator algebras. Cham: Springer Open, 2017.

\bibitem{Oml} M. Omladi\v c{}, On operators preserving commutativity, J. Funct. Anal. {\bf 66}, (1986), 105--122.

\bibitem{OS} M. Omladi\v c{} and P. \v Semrl, Preserving diagonalisability, Linear Algebra Appl. {\bf 285} (1998), no.~1-3, 165--179.

\bibitem{UHL} U. Uhlhorn, Representation of symmetry transformations in quantum mechanics, Ark. Fysik
{\bf 23}, 307–340 (1963).

%\bibitem{Weyl} H. Weyl, {\"U}ber beschr{\"a}nkte quadratische Formen, deren Differenz vollstetig ist., Rend. Circ. Mat. Palermo 27, 373--392, 402 (1909). 
\end{thebibliography}

\end{document}

